% GENERATED FILE. Edit canonical lessons, then run npm run generate:books.
% canonical-source-hash: fnv1a64:1f9ae04f65b43c34
% canonical-lessons: FR-C401-accelerer, FR-C401-braquer, FR-C401-deraper, FR-C401-caler, FR-C401-pedale

\chapter{Au volant}
\label{ch:fr-au-volant}
\hlchaptermodality{401}

\begin{chapteropening}
\textbf{By the end of this chapter:} \emph{I can describe how someone drives: speeding up, turning the wheel, skidding, stalling and the pedals.}

Complete the last lesson of chapter 401: I can describe how someone drives: speeding up, turning the wheel, skidding, stalling and the pedals.
\end{chapteropening}

\section[accélérer]{accélérer — to speed up, to accelerate}

\label{lesson:FR-C401-accelerer}



\begin{quote}
Before the new one: say the French for the slip road, then the French for to turn off.
\end{quote}

\subsection*{You'll want to know: accélérer}
\textbf{accélérer} — \textquotedblleft{}to speed up, to accelerate\textquotedblright{}.

\begin{quote}
Sur l'autoroute, mon père accélère un peu trop.
\end{quote}

\begin{quote}
\emph{On the motorway my father speeds up a little too much.}
\end{quote}

\emph{L'accélérateur} is the accelerator pedal.

\begin{cousinweb}
\emph{accélérer} is a learned borrowing from Latin \emph{accelerāre} 'to hasten', built on \emph{celer} 'swift' (Lewis \& Short).
\end{cousinweb}

\begin{grammarlens}[title={What you've built}]
One more word for telling, explaining and describing at B1.
\end{grammarlens}

\subsection*{Your turn}
\begin{itemize}
\raggedright
  \item \emph{Say it:} \emph{accélérer}
  \item \emph{Say it:} \emph{accélérer}, once more
  \item \emph{Say it:} a sentence of your own with \emph{accélérer}
  \item \emph{Recall:} say the French for the slip road, then the French for to turn off, then say \emph{accélérer} again
\end{itemize}

\begin{tcolorbox}[breakable,colback=teal!4,colframe=teal!35!black,title={Before you move on}]
What does accélérer mean? (\textquotedblleft{}To speed up, to accelerate\textquotedblright{}.) Say it once more.
\end{tcolorbox}
\section[braquer]{braquer — to turn the wheel (hard); to aim}

\label{lesson:FR-C401-braquer}



\begin{quote}
Before the new one: say the French for to turn off, then the French for to speed up.
\end{quote}

\subsection*{You'll want to know: braquer}
\textbf{braquer} — \textquotedblleft{}to turn the wheel (hard); to aim\textquotedblright{}.

\begin{quote}
Elle a braqué à gauche pour éviter le chat.
\end{quote}

\begin{quote}
\emph{She turned the wheel hard to the left to avoid the cat.}
\end{quote}

\emph{Braquer une banque} is slang for to hold up a bank.

\begin{cousinweb}
The origin of \emph{braquer} is not fully settled. It is usually traced to a Vulgar Latin verb built on Latin \emph{bracchium} 'arm' (Lewis \& Short): first to point or turn something with the arm.
\end{cousinweb}

\begin{grammarlens}[title={What you've built}]
One more word for telling, explaining and describing at B1.
\end{grammarlens}

\subsection*{Your turn}
\begin{itemize}
\raggedright
  \item \emph{Say it:} \emph{braquer}
  \item \emph{Say it:} \emph{braquer}, once more
  \item \emph{Say it:} a sentence of your own with \emph{braquer}
  \item \emph{Recall:} say the French for to turn off, then the French for to speed up, then say \emph{braquer} again
\end{itemize}

\begin{tcolorbox}[breakable,colback=teal!4,colframe=teal!35!black,title={Before you move on}]
What does braquer mean? (\textquotedblleft{}To turn the wheel (hard); to aim\textquotedblright{}.) Say it once more.
\end{tcolorbox}
\section[déraper]{déraper — to skid}

\label{lesson:FR-C401-deraper}



\begin{quote}
Before the new one: say the French for to speed up, then the French for to turn the wheel.
\end{quote}

\subsection*{You'll want to know: déraper}
\textbf{déraper} — \textquotedblleft{}to skid\textquotedblright{}.

\begin{quote}
Avec la pluie, la voiture a dérapé dans la descente.
\end{quote}

\begin{quote}
\emph{Because of the rain the car skidded on the way down.}
\end{quote}

Also figurative: \emph{la conversation a dérapé} — the talk got out of hand.

\begin{cousinweb}
\emph{déraper} is usually traced to Provençal \emph{derapar}, first said of an anchor pulled loose from the sea bed, built on a Germanic verb for seizing. EtymDB links it with Germanic cousins such as Old Norse \emph{hreppa} 'to grasp' (EtymDB).
\end{cousinweb}

\begin{grammarlens}[title={What you've built}]
One more word for telling, explaining and describing at B1.
\end{grammarlens}

\subsection*{Your turn}
\begin{itemize}
\raggedright
  \item \emph{Say it:} \emph{déraper}
  \item \emph{Say it:} \emph{déraper}, once more
  \item \emph{Say it:} a sentence of your own with \emph{déraper}
  \item \emph{Recall:} say the French for to speed up, then the French for to turn the wheel, then say \emph{déraper} again
\end{itemize}

\begin{tcolorbox}[breakable,colback=teal!4,colframe=teal!35!black,title={Before you move on}]
What does déraper mean? (\textquotedblleft{}To skid\textquotedblright{}.) Say it once more.
\end{tcolorbox}
\section[caler]{caler — to stall (an engine); to wedge}

\label{lesson:FR-C401-caler}



\begin{quote}
Before the new one: say the French for to turn the wheel, then the French for to skid.
\end{quote}

\subsection*{You'll want to know: caler}
\textbf{caler} — \textquotedblleft{}to stall (an engine); to wedge\textquotedblright{}.

\begin{quote}
Au feu rouge, j'ai calé deux fois.
\end{quote}

\begin{quote}
\emph{At the red light I stalled twice.}
\end{quote}

\emph{Je cale} also means I give up, I cannot eat any more.

\begin{cousinweb}
\emph{caler} in the sense 'to stall' was borrowed from Provençal \emph{calar} 'to lower a sail', which goes back to Greek \emph{khalân} 'to slacken' (EtymDB). A boat that lowers its sail stops, and so does an engine.
\end{cousinweb}

\begin{grammarlens}[title={What you've built}]
One more word for telling, explaining and describing at B1.
\end{grammarlens}

\subsection*{Your turn}
\begin{itemize}
\raggedright
  \item \emph{Say it:} \emph{caler}
  \item \emph{Say it:} \emph{caler}, once more
  \item \emph{Say it:} a sentence of your own with \emph{caler}
  \item \emph{Recall:} say the French for to turn the wheel, then the French for to skid, then say \emph{caler} again
\end{itemize}

\begin{tcolorbox}[breakable,colback=teal!4,colframe=teal!35!black,title={Before you move on}]
What does caler mean? (\textquotedblleft{}To stall (an engine); to wedge\textquotedblright{}.) Say it once more.
\end{tcolorbox}
\section[la pédale]{la pédale — the pedal}

\label{lesson:FR-C401-pedale}



\begin{quote}
Before the new one: say the French for to skid, then the French for to stall.
\end{quote}

\subsection*{You'll want to know: la pédale}
\textbf{la pédale} — \textquotedblleft{}the pedal\textquotedblright{}.

\begin{quote}
Le pied droit reste près de la pédale.
\end{quote}

\begin{quote}
\emph{The right foot stays near the pedal.}
\end{quote}

\emph{Pédaler} — to pedal, on a bike.

\begin{cousinweb}
\emph{pédale} was borrowed from Italian \emph{pedale}, which is inherited from Latin \emph{pedālis} 'of the foot', from \emph{pēs} 'foot' (EtymDB).
\end{cousinweb}

\begin{grammarlens}[title={What you've built}]
One more word for telling, explaining and describing at B1.
\end{grammarlens}

\subsection*{Your turn}
\begin{itemize}
\raggedright
  \item \emph{Say it:} \emph{la pédale}
  \item \emph{Say it:} \emph{la pédale}, once more
  \item \emph{Say it:} a sentence of your own with \emph{la pédale}
  \item \emph{Recall:} say the French for to skid, then the French for to stall, then say \emph{la pédale} again
\end{itemize}

\begin{tcolorbox}[breakable,colback=teal!4,colframe=teal!35!black,title={Before you move on}]
What does la pédale mean? (\textquotedblleft{}The pedal\textquotedblright{}.) Say it once more.
\end{tcolorbox}
